\documentclass[12pt]{article}
\usepackage{amsmath,amssymb}

\begin{document}

\section*{An Extremal Problem from Graph Theory}

The problem under consideration is the following. Let $n$ and $k$ be given natural numbers ($2 \le k \le n$). How many edges can a graph with $n$ vertices contain, in which each pair of vertices is connected by at most one edge, if it contains no complete subgraph with $k$ vertices? (A graph is said to be complete if every two of its vertices are connected by exactly one edge.)

We first describe a graph $D(n,k)$ corresponding to the given pair $(n,k)$ of natural numbers, which---in the elegant formulation of Prof.\ D.~Kőnig---can be defined as follows: we number $n$ vertices in order by $1,2,\dots,n$ and connect two vertices if and only if the numbers assigned to them are incongruent modulo $(k-1)$. It will be proved that the graph $D(n,k)$ thus obtained (for $n=7$, $k=4$, see the figure) provides the solution, and indeed the unique solution, of the above extremal problem.

The number of edges of this graph is:
\[
\frac{k-2}{2(k-1)}\,(n^2 - r^2) + \binom{r}{2},
\]
where $r$ is defined by
\[
n \equiv r \pmod{k-1}, \qquad 0 \le r < k-1.
\]

[For $k=3$, this theorem was already found in 1907 by W.~Mantel (see \emph{Fortschritte d.\ Math.}, B.\ 38, p.\ 270).]

The problem solved here can also be formulated---in purely combinatorial terms---as follows. One is to determine a smallest possible set $M$ of combinations of the second class (without repetition) from $n$ elements such that every combination of the $k$th class (without repetition) of these $n$ elements contains at least one element of $M$.

For infinite graphs the following theorem will be proved. A graph with countably infinite vertices always contains an infinite complete subgraph if the following condition is satisfied for a natural number $k$: whenever $k$ vertices of the graph are chosen arbitrarily, there are always at least two among them that are connected by an edge of the graph.

For $k=3$, this theorem is due to Mr.\ G.~Szekeres, who was led to such graphs through group-theoretic investigations connected with the Burnside problem.

\bigskip
\begin{flushright}
Paul Turán
\end{flushright}

\end{document}

